%% LyX 2.4.4 created this file.  For more info, see https://www.lyx.org/.
%% Do not edit unless you really know what you are doing.
\documentclass[english,t]{beamer}
\usepackage{mathptmx}
\usepackage[T1]{fontenc}
\usepackage[latin9]{inputenc}
\usepackage{amssymb}
\usepackage{cancel}
\PassOptionsToPackage{normalem}{ulem}
\usepackage{ulem}

\makeatletter

%%%%%%%%%%%%%%%%%%%%%%%%%%%%%% LyX specific LaTeX commands.
%% Because html converters don't know tabularnewline
\providecommand{\tabularnewline}{\\}

%%%%%%%%%%%%%%%%%%%%%%%%%%%%%% Textclass specific LaTeX commands.
% this default might be overridden by plain title style
\newcommand\makebeamertitle{\frame{\maketitle}}%
% (ERT) argument for the TOC
\AtBeginDocument{%
  \let\origtableofcontents=\tableofcontents
  \def\tableofcontents{\@ifnextchar[{\origtableofcontents}{\gobbletableofcontents}}
  \def\gobbletableofcontents#1{\origtableofcontents}
}

%%%%%%%%%%%%%%%%%%%%%%%%%%%%%% User specified LaTeX commands.
\usetheme{Warsaw}
% or ...

\setbeamercovered{transparent}
% or whatever (possibly just delete it)
\setbeamercovered{invisible} %makes overlays invisible


%gets rid of navigation symbols
\setbeamertemplate{navigation symbols}{}

\usepackage{multicol}

\usepackage{tikz}
\usetikzlibrary{calc}
\usetikzlibrary{plotmarks}
\usepackage{pgfplots}
\pgfplotsset{compat=newest}

\usepackage{calc}
\usepackage[nomessages]{fp}

\usepackage{advdate}    % Advancing/saving dates
\usepackage{datetime}   % Dates formatting
\usepackage{datenumber} % Counters for dates
\usepackage[super]{nth}

\usepackage{etex}


%\usepackage{pgfpages}
%\pgfpagesuselayout{4 on 1}[a4paper, landscape, border shrink=7mm]
%\pgfpageslogicalpageoptions{1}{border code=\pgfusepath{stroke}}
%\pgfpageslogicalpageoptions{2}{border code=\pgfusepath{stroke}}
%\pgfpageslogicalpageoptions{3}{border code=\pgfusepath{stroke}}
%\pgfpageslogicalpageoptions{4}{border code=\pgfusepath{stroke}}
%%%Don't forget to add 'handout'

\makeatother

\usepackage{babel}
\begin{document}
\newcommand{\xmax}{2000}
\newcommand{\xmin}{0}
\newcommand{\xstep}{0} %must be xmin+n
\newcommand{\ymax}{500}
\newcommand{\ymin}{0}
\newcommand{\ystep}{0} %must be ymin+n

\newcounter{countEx} \setcounter{countEx}{0}
\newcounter{countProb} \setcounter{countProb}{0}
\resetcounteronoverlays{countEx} \resetcounteronoverlays{countProb}

\newcommand{\ex}[3]{
	\addtocounter{countEx}{1}
	\only<#1-#2>{\begin{exampleblock} {Example \chNum.\secNum.\arabic{countEx}.} #3	\end{exampleblock}}
}

\newcommand{\prob}[4]{
	\addtocounter{countProb}{1}
	\only<#1-#2>{\begin{exampleblock} {Problem  \chNum.\secNum.\arabic{countProb}.\,#3} #4	\end{exampleblock}}
}

\newcommand{\yline}[4]{
	(#2 - #4)/(#1 - #3)*(x - #1) + #2
}

\beamerdefaultoverlayspecification{<+->}

%%%%Chapter and Section numbers%%%%%
\newcommand{\chNum}{7}
\newcommand{\secNum}{1}
\newcommand{\secTitle}{Sets}

\title[Section \chNum.\secNum: \secTitle]{Section \chNum.\secNum: \secTitle}

\subtitle{Finite Mathematics~~MTH 132~~Spring 2014}

\author{Dr.\,James Ashe}

%\titlegraphic{\includegraphics[height=1.5cm]{/home/jim/JamesAshe/JCSU/images/jcsu_bull}}

\institute{Johnson C. Smith University}

\date{{}}

\makebeamertitle 
\begin{frame}{Chapter 7: Sets and Probability}

\begin{columns}


\column{11cm}
\begin{itemize}
\item []Chapter $7$: Sets and Probability

\begin{itemize}
\item [7.1:] Sets
\item [7.2:] Venn Diagrams
\item [7.3:] Introduction to Probability 
\end{itemize}
\end{itemize}
\end{columns}
\end{frame}

\begin{frame}{Sets}

\vspace{-.75cm}
\begin{columns}


\column{5.5cm}
\begin{block}{Set}
 A \emph{set }is a well defined collection of distinct objects.

\begin{itemize}
\item Usually written as capital letters, e.g., \emph{A, B, C, }etc.
\end{itemize}
\end{block}
\footnotesize
\begin{itemize}
\item <2->[ex.] $\left\{ 0,2,4\right\} $
\item <3->[ex.] $N=\left\{ Billy,\,Bob,\,Bubba\right\} $
\item <4->[ex.] The set of letters in the alphabet

\begin{itemize}
\item <5->[]$\{a,b,c,\dots,x,y,z\}$
\end{itemize}
\item <6->[ex.] The set of positive integers.

\begin{itemize}
\item <7->[]$\left\{ 1,2,3,4,\dots\right\} $
\end{itemize}
\item <8->[ex.] The set of positive integers$<0$.

\begin{itemize}
\item <9->[]$\emptyset$ ($\emptyset\neq0$)
\end{itemize}
\end{itemize}

\column{5.5cm}
\begin{block}<11->{Element}
 The objects in the set are called \emph{elements} or \emph{members}.

\begin{itemize}
\item $x\in A$ means $x$ is an element in $A$.
\item $x\notin A$ means $x$ is not an element in $A$.
\end{itemize}
\end{block}
\footnotesize
\begin{itemize}
\item <12->[ex.] $2\in\left\{ 0,2,4\right\} $
\item <13->[ex.] $Carlton\notin N$
\item <14->[ex.] $h\in\{a,b,c,\dots,x,y,z\}$
\item <15->[ex.] $3.14\notin\left\{ 1,2,3,4,\dots\right\} $
\item <16->[ex.] $x\notin\emptyset$
\end{itemize}
\end{columns}
\end{frame}

\begin{frame}{Sets}

\vspace{-.75cm}
\begin{columns}


\column{4.5cm}
\begin{itemize}
\item <1->Two sets are \emph{equal} if the contain the same elements.

\begin{itemize}
\item <3->Order does not matter.
\item <5->We do not repeat elements
\end{itemize}
\item <7->[]A set can be given by,

\begin{itemize}
\item <8->[1.] description 
\item <10->[2.] \emph{tabular }form
\item <12->[3.] \emph{set-builder} notation
\end{itemize}
\end{itemize}

\column{6.5cm}

\footnotesize
\begin{itemize}
\item []
\item []
\item <2->[ex.]$\{0,2,4\}\neq\left\{ 0,2,4,1\right\} $
\item <4->[ex.]$\{0,2,4\}=\{2,4,0\}=\{4,0,2\}$
\item <6->[]Do not write $\{0,2,2,4\}$
\item []
\item []
\item <9->[ex.] The set of positive integers
\item <11->[ex.] $\{1,2,3,4,\dots\}$
\item <13->[ex.] $\left\{ x\mid x>0\mbox{ and }x\mbox{ is an integer}\right\} $
\end{itemize}
\end{columns}
\prob{14}{14}{Write the following set in \emph{tabular} and \emph{set-builder} form.}{The set of even, positive integers.}

\prob{15}{15}{Decide whether the following are \emph{true} or \emph{false}.}{$\{1,2,3,4\}=\{4,3,1,2\}$}\prob{16}{16}{Decide whether the following are \emph{true} or \emph{false}.}{$\{a,b,\dots,z\}=\{b,c,\dots,z\}$}
\end{frame}

\begin{frame}{Subsets}

\vspace{-.75cm}
\begin{columns}


\column{5.5cm}
\begin{block}{Subset}
 a set $A$ is a \emph{subset} of a set $B$ if every element in
$A$ is also in $B$.

\begin{itemize}
\item []<2->$A\subseteq B$: $A$ is a subset of $B$
\item []<3->$A\nsubseteq B$: $A$ is \uline{not} a subset of $B$
\end{itemize}
\end{block}

\footnotesize
\begin{itemize}
\item <4->[ex.]

\begin{itemize}
\item <4->[]%
\begin{tabular}{l}
$A=\left\{ {\color{blue}3,4,5,6}\right\} $\tabularnewline
$B=\left\{ 2,{\color{blue}3,4,5,6}\mathpunct{\color{blue},}7,8\right\} $\tabularnewline
$C=\left\{ {\color{blue}3,4,5,6},9\right\} $\tabularnewline
\end{tabular}
\end{itemize}
\item <5-> $A\subseteq B$
\item <6->$C\nsubseteq B$
\item <7->$\emptyset\subseteq A$, $\emptyset\subseteq B$, and $\emptyset\subseteq C$
\item <8->$B\subseteq B$, $C\subseteq C$, $A\subseteq A$
\end{itemize}

\column{5.5cm}
\begin{block}<12->{Proper Subset}
 a set $A$ is a \emph{proper subset} of a set $B$ if every element
in $A$ is also in $B$, and $A\neq B$.

\begin{itemize}
\item []$A\subset B$: $A$ is a proper subset of $B$.
\end{itemize}
\end{block}
\begin{itemize}
\item <13->[ex.]

\begin{itemize}
\item <13->[]%
\begin{tabular}{l}
$A=\left\{ {\color{blue}3,4,5,6}\right\} $\tabularnewline
$B=\left\{ 2,{\color{blue}3,4,5,6}\mathpunct{\color{blue},}7,8\right\} $\tabularnewline
$C=\left\{ {\color{blue}3,4,5,6},9\right\} $\tabularnewline
\end{tabular}
\end{itemize}
\item <14->$A\subset B$
\item <15->$B$ is not a proper subset of $B$. 
\end{itemize}
\end{columns}
\end{frame}

\begin{frame}{Subsets}

\begin{block}<1->{}

\begin{itemize}
\item $A\subseteq A$ for any set $A$.
\item $\emptyset\subseteq A$
\item A set with $k$ elements has $2^{k}$ subsets.
\end{itemize}
\end{block}
\begin{itemize}
\item <2->[ex.]$\{1,2\}$ has $2^{2}$ subsets.

\begin{itemize}
\item <3->[]$\{1,2\}$, $\{1\}$, $\{2\}$, $\emptyset$.
\end{itemize}
\end{itemize}
\prob{4}{4}{Let $A=\{2,4,6,10\}$, $B=\{2,10\}$, and $C=\{2,4,6,10\}$. Decide whether the following are true or false.}{$B \subset A$}

\prob{5}{5}{Let $A=\{2,4,6,10\}$, $B=\{2,10\}$, and $C=\{2,4,6,10\}$. Decide whether the following are true or false.}{$A \subseteq B$}

\prob{6}{6}{Let $A=\{2,4,6,10\}$, $B=\{2,10\}$, and $C=\{2,4,6,10\}$. Decide whether the following are true or false.}{$A \subset C$}

\prob{7}{7}{Let $A=\{2,4,6,10\}$, $B=\{2,10\}$, and $C=\{2,4,6,10\}$. Decide whether the following are true or false.}{$A \subseteq C$ and $C \subseteq A$}

\prob{8}{8}{Let $A=\{2,4,6,10\}$, $B=\{2,10\}$, and $C=\{2,4,6,10\}$. Decide whether the following are true or false.}{$A=C$}

\prob{9}{9}{Let $E=\{2,4,6\}$}{How many subsets does $E$ have? List them.}
\end{frame}

\end{document}
